% --- Archon physics formalization source begin ---
% archon:physics
% archon:covers PhyXMiniProblems/problem_phyx_mini_0587.lean
% archon:source-report reports/phyx_mini/problem_phyx_mini_0587.source.json

\chapter{Physics problem phyx\_mini\_0587}
\label{ch:PhyXMiniProblems_problem_phyx_mini_0587}

\paragraph{Problem source.}
\#\# Physical scenario

An electron is trapped in a one-dimensional infinite well and is in its first excited state. Figure indicates the five longest wavelengths of light that the electron could absorb in transitions from this initial state via a single photon absorption: \textbackslash{}( \textbackslash{}lambda\_a = 80.78\textbackslash{},\textbackslash{}text\{nm\} \textbackslash{}), \textbackslash{}( \textbackslash{}lambda\_b = 33.66\textbackslash{},\textbackslash{}text\{nm\} \textbackslash{}), \textbackslash{}( \textbackslash{}lambda\_c = 19.23\textbackslash{},\textbackslash{}text\{nm\} \textbackslash{}), \textbackslash{}( \textbackslash{}lambda\_d = 12.62\textbackslash{},\textbackslash{}text\{nm\} \textbackslash{}), and \textbackslash{}( \textbackslash{}lambda\_e = 8.98\textbackslash{},\textbackslash{}text\{nm\} \textbackslash{}).

\#\# Question

What is the width of the potential well?

\#\# Answer choices

- A: A: 280pm

- B: B: 700pm

- C: C: 350pm

- D: D: 175pm

\#\# Figure caption (auxiliary; use the image as primary evidence)

The image is a linear diagram along a horizontal axis labeled "λ (nm)" with a starting point marked as 0. The axis is divided into segments, each representing an unspecified numerical scale.

Five distinct labeled points with red dots are plotted along the axis. These points are labeled as follows from left to right:

1. λₑ
2. λ\_d
3. λ\_c
4. λ\_b
5. λ\_a

Each label is positioned slightly above the corresponding red dot. The placement of the dots suggests a sequential ordering but the exact numerical values or scale are not provided.

\#\# Dataset metadata

Quantum Phenomena; Physical Model Grounding Reasoning, Multi-Formula Reasoning

\paragraph{Recorded answer/context.}
C: C: 350pm

\paragraph{Figure/image path.}
/root/proposal\_for\_physic/science-mango/phyx\_mini\_run/phyx\_data/test\_image/587.png

\paragraph{Formalization target.}
create a compiling Lean file with sorry bodies at `PhyXMiniProblems/problem\_phyx\_mini\_0587.lean`. The Lean declarations must preserve the physical quantities, dimensions or dimensional roles, figure labels, governing-law hypotheses, and final relation expressed by this problem.
Use Mathlib/Physlib names found through LeanExplore where available. If a physics API is missing, introduce faithful local abstractions rather than scalar placeholder aliases.

\begin{theorem}[Physics formalization target]
\label{thm:physics:phyx_mini_0587:target}
\lean{PhyXMiniProblems.ProblemPhyXMini0587.problem_phyx_mini_0587}
\uses{def:physics:phyx-mini-0587:phyxminiproblems-problemphyxmini0587-infinitewellabsorptionsetup, def:physics:phyx-mini-0587:phyxminiproblems-problemphyxmini0587-matchesfirstexcitedelectroninfinitewell, def:physics:phyx-mini-0587:phyxminiproblems-problemphyxmini0587-matchesgivenabsorptionwavelengthreadouts, def:physics:phyx-mini-0587:phyxminiproblems-problemphyxmini0587-matchessuppliedabsorptionfigure, def:physics:phyx-mini-0587:phyxminiproblems-problemphyxmini0587-matchesfivelongestabsorptiontransitions, def:physics:phyx-mini-0587:phyxminiproblems-problemphyxmini0587-usesstandardreferencedata, def:physics:phyx-mini-0587:phyxminiproblems-problemphyxmini0587-hasphysicalinfinitewellparameters, def:physics:phyx-mini-0587:phyxminiproblems-problemphyxmini0587-satisfiesinfinitesquarewellspectrum, def:physics:phyx-mini-0587:phyxminiproblems-problemphyxmini0587-satisfiessinglephotonabsorptionlaw, def:physics:phyx-mini-0587:phyxminiproblems-problemphyxmini0587-agreeswhenroundedtonearestpicometer, def:physics:phyx-mini-0587:phyxminiproblems-problemphyxmini0587-isuniquenearestwellwidthchoice}
The assigned autoformalize agent should translate this physics problem into Lean declarations in the covered file, with theorem and lemma proof bodies written as `by sorry`.
\end{theorem}
\begin{proof}
This is an autoformalization task, not a proof task. Produce statements that can later be proved by the physics prover without weakening the physical model.
\end{proof}

% --- Archon declaration topology begin ---
\section{Declaration topology}
\begin{definition}[Length Quantity]
  \label{def:physics:phyx-mini-0587:phyxminiproblems-problemphyxmini0587-lengthquantity}
  \lean{PhyXMiniProblems.ProblemPhyXMini0587.LengthQuantity}
  A nonnegative, unit-independent physical length
\end{definition}
\begin{proof}
  This is the declared model definition of Length Quantity.
\end{proof}
\begin{definition}[Mass Quantity]
  \label{def:physics:phyx-mini-0587:phyxminiproblems-problemphyxmini0587-massquantity}
  \lean{PhyXMiniProblems.ProblemPhyXMini0587.MassQuantity}
  A nonnegative, unit-independent physical mass
\end{definition}
\begin{proof}
  This is the declared model definition of Mass Quantity.
\end{proof}
\begin{definition}[action Dimension]
  \label{def:physics:phyx-mini-0587:phyxminiproblems-problemphyxmini0587-actiondimension}
  \lean{PhyXMiniProblems.ProblemPhyXMini0587.actionDimension}
  The physical dimension of action, mass * length\textasciicircum{}2 / time
\end{definition}
\begin{proof}
  This is the declared model definition of action Dimension.
\end{proof}
\begin{definition}[Action Quantity]
  \label{def:physics:phyx-mini-0587:phyxminiproblems-problemphyxmini0587-actionquantity}
  \lean{PhyXMiniProblems.ProblemPhyXMini0587.ActionQuantity}
  \uses{def:physics:phyx-mini-0587:phyxminiproblems-problemphyxmini0587-actiondimension}
  A nonnegative, unit-independent physical action
\end{definition}
\begin{proof}
  This is the declared model definition of Action Quantity.
\end{proof}
\begin{definition}[Light Speed Quantity]
  \label{def:physics:phyx-mini-0587:phyxminiproblems-problemphyxmini0587-lightspeedquantity}
  \lean{PhyXMiniProblems.ProblemPhyXMini0587.LightSpeedQuantity}
  A signed, unit-independent physical speed, matching Physlib's type for c
\end{definition}
\begin{proof}
  This is the declared model definition of Light Speed Quantity.
\end{proof}
\begin{definition}[length Readout]
  \label{def:physics:phyx-mini-0587:phyxminiproblems-problemphyxmini0587-lengthreadout}
  \lean{PhyXMiniProblems.ProblemPhyXMini0587.lengthReadout}
  \uses{def:physics:phyx-mini-0587:phyxminiproblems-problemphyxmini0587-lengthquantity}
  Read a physical length in a selected length unit
\end{definition}
\begin{proof}
  This is the declared model definition of length Readout.
\end{proof}
\begin{definition}[length In Meters]
  \label{def:physics:phyx-mini-0587:phyxminiproblems-problemphyxmini0587-lengthinmeters}
  \lean{PhyXMiniProblems.ProblemPhyXMini0587.lengthInMeters}
  \uses{def:physics:phyx-mini-0587:phyxminiproblems-problemphyxmini0587-lengthquantity, def:physics:phyx-mini-0587:phyxminiproblems-problemphyxmini0587-lengthreadout}
  Read a physical length in metres
\end{definition}
\begin{proof}
  This is the declared model definition of length In Meters.
\end{proof}
\begin{definition}[length In Nanometers]
  \label{def:physics:phyx-mini-0587:phyxminiproblems-problemphyxmini0587-lengthinnanometers}
  \lean{PhyXMiniProblems.ProblemPhyXMini0587.lengthInNanometers}
  \uses{def:physics:phyx-mini-0587:phyxminiproblems-problemphyxmini0587-lengthquantity, def:physics:phyx-mini-0587:phyxminiproblems-problemphyxmini0587-lengthreadout}
  Read a physical length in nanometres
\end{definition}
\begin{proof}
  This is the declared model definition of length In Nanometers.
\end{proof}
\begin{definition}[length In Picometers]
  \label{def:physics:phyx-mini-0587:phyxminiproblems-problemphyxmini0587-lengthinpicometers}
  \lean{PhyXMiniProblems.ProblemPhyXMini0587.lengthInPicometers}
  \uses{def:physics:phyx-mini-0587:phyxminiproblems-problemphyxmini0587-lengthquantity, def:physics:phyx-mini-0587:phyxminiproblems-problemphyxmini0587-lengthreadout}
  Read a physical length in picometres
\end{definition}
\begin{proof}
  This is the declared model definition of length In Picometers.
\end{proof}
\begin{definition}[mass In Kilograms]
  \label{def:physics:phyx-mini-0587:phyxminiproblems-problemphyxmini0587-massinkilograms}
  \lean{PhyXMiniProblems.ProblemPhyXMini0587.massInKilograms}
  \uses{def:physics:phyx-mini-0587:phyxminiproblems-problemphyxmini0587-massquantity}
  Read a mass in kilograms
\end{definition}
\begin{proof}
  This is the declared model definition of mass In Kilograms.
\end{proof}
\begin{definition}[action In Joule Seconds]
  \label{def:physics:phyx-mini-0587:phyxminiproblems-problemphyxmini0587-actioninjouleseconds}
  \lean{PhyXMiniProblems.ProblemPhyXMini0587.actionInJouleSeconds}
  \uses{def:physics:phyx-mini-0587:phyxminiproblems-problemphyxmini0587-actionquantity}
  Read an action in coherent SI units, joule-seconds
\end{definition}
\begin{proof}
  This is the declared model definition of action In Joule Seconds.
\end{proof}
\begin{definition}[speed In Meters Per Second]
  \label{def:physics:phyx-mini-0587:phyxminiproblems-problemphyxmini0587-speedinmeterspersecond}
  \lean{PhyXMiniProblems.ProblemPhyXMini0587.speedInMetersPerSecond}
  \uses{def:physics:phyx-mini-0587:phyxminiproblems-problemphyxmini0587-lightspeedquantity}
  Read a speed in coherent SI units, metres per second
\end{definition}
\begin{proof}
  This is the declared model definition of speed In Meters Per Second.
\end{proof}
\begin{definition}[energy In Joules]
  \label{def:physics:phyx-mini-0587:phyxminiproblems-problemphyxmini0587-energyinjoules}
  \lean{PhyXMiniProblems.ProblemPhyXMini0587.energyInJoules}
  Coherent-SI readout of an energy, in joules
\end{definition}
\begin{proof}
  This is the declared model definition of energy In Joules.
\end{proof}
\begin{definition}[Particle Species]
  \label{def:physics:phyx-mini-0587:phyxminiproblems-problemphyxmini0587-particlespecies}
  \lean{PhyXMiniProblems.ProblemPhyXMini0587.ParticleSpecies}
  Particle species distinguished by the problem statement
\end{definition}
\begin{proof}
  This is the declared model definition of Particle Species.
\end{proof}
\begin{definition}[Potential Well Model]
  \label{def:physics:phyx-mini-0587:phyxminiproblems-problemphyxmini0587-potentialwellmodel}
  \lean{PhyXMiniProblems.ProblemPhyXMini0587.PotentialWellModel}
  The confinement model used for the electron
\end{definition}
\begin{proof}
  This is the declared model definition of Potential Well Model.
\end{proof}
\begin{definition}[Absorption Line Label]
  \label{def:physics:phyx-mini-0587:phyxminiproblems-problemphyxmini0587-absorptionlinelabel}
  \lean{PhyXMiniProblems.ProblemPhyXMini0587.AbsorptionLineLabel}
  Labels attached to the five plotted absorption wavelengths
\end{definition}
\begin{proof}
  This is the declared model definition of Absorption Line Label.
\end{proof}
\begin{definition}[Figure Axis Role]
  \label{def:physics:phyx-mini-0587:phyxminiproblems-problemphyxmini0587-figureaxisrole}
  \lean{PhyXMiniProblems.ProblemPhyXMini0587.FigureAxisRole}
  Physical role of the horizontal axis in the supplied diagram
\end{definition}
\begin{proof}
  This is the declared model definition of Figure Axis Role.
\end{proof}
\begin{definition}[Absorption Line]
  \label{def:physics:phyx-mini-0587:phyxminiproblems-problemphyxmini0587-absorptionline}
  \lean{PhyXMiniProblems.ProblemPhyXMini0587.AbsorptionLine}
  \uses{def:physics:phyx-mini-0587:phyxminiproblems-problemphyxmini0587-lengthquantity}
  Data carried by one single-photon absorption line
\end{definition}
\begin{proof}
  This is the declared model definition of Absorption Line.
\end{proof}
\begin{definition}[Infinite Well Absorption Figure]
  \label{def:physics:phyx-mini-0587:phyxminiproblems-problemphyxmini0587-infinitewellabsorptionfigure}
  \lean{PhyXMiniProblems.ProblemPhyXMini0587.InfiniteWellAbsorptionFigure}
  \uses{def:physics:phyx-mini-0587:phyxminiproblems-problemphyxmini0587-absorptionlinelabel, def:physics:phyx-mini-0587:phyxminiproblems-problemphyxmini0587-figureaxisrole}
  Literal qualitative content visible in image 587.png
\end{definition}
\begin{proof}
  This is the declared model definition of Infinite Well Absorption Figure.
\end{proof}
\begin{definition}[Infinite Well Absorption Setup]
  \label{def:physics:phyx-mini-0587:phyxminiproblems-problemphyxmini0587-infinitewellabsorptionsetup}
  \lean{PhyXMiniProblems.ProblemPhyXMini0587.InfiniteWellAbsorptionSetup}
  \uses{def:physics:phyx-mini-0587:phyxminiproblems-problemphyxmini0587-lengthquantity, def:physics:phyx-mini-0587:phyxminiproblems-problemphyxmini0587-massquantity, def:physics:phyx-mini-0587:phyxminiproblems-problemphyxmini0587-actionquantity, def:physics:phyx-mini-0587:phyxminiproblems-problemphyxmini0587-lightspeedquantity, def:physics:phyx-mini-0587:phyxminiproblems-problemphyxmini0587-particlespecies, def:physics:phyx-mini-0587:phyxminiproblems-problemphyxmini0587-potentialwellmodel, def:physics:phyx-mini-0587:phyxminiproblems-problemphyxmini0587-absorptionlinelabel, def:physics:phyx-mini-0587:phyxminiproblems-problemphyxmini0587-absorptionline, def:physics:phyx-mini-0587:phyxminiproblems-problemphyxmini0587-infinitewellabsorptionfigure}
  Independent physical quantities in the absorption experiment. In particular, wellWidth is stored independently of the displayed answer choices
\end{definition}
\begin{proof}
  This is the declared model definition of Infinite Well Absorption Setup.
\end{proof}
\begin{definition}[Matches First Excited Electron Infinite Well]
  \label{def:physics:phyx-mini-0587:phyxminiproblems-problemphyxmini0587-matchesfirstexcitedelectroninfinitewell}
  \lean{PhyXMiniProblems.ProblemPhyXMini0587.MatchesFirstExcitedElectronInfiniteWell}
  \uses{def:physics:phyx-mini-0587:phyxminiproblems-problemphyxmini0587-infinitewellabsorptionsetup}
  The prose scenario: an electron in the first excited state of a 1D infinite well
\end{definition}
\begin{proof}
  This is the declared model definition of Matches First Excited Electron Infinite Well.
\end{proof}
\begin{definition}[Agrees With Two Decimal Nanometer Readout]
  \label{def:physics:phyx-mini-0587:phyxminiproblems-problemphyxmini0587-agreeswithtwodecimalnanometerreadout}
  \lean{PhyXMiniProblems.ProblemPhyXMini0587.AgreesWithTwoDecimalNanometerReadout}
  \uses{def:physics:phyx-mini-0587:phyxminiproblems-problemphyxmini0587-lengthquantity, def:physics:phyx-mini-0587:phyxminiproblems-problemphyxmini0587-lengthinnanometers}
  Agreement with a wavelength printed to two decimal places in nanometres
\end{definition}
\begin{proof}
  This is the declared model definition of Agrees With Two Decimal Nanometer Readout.
\end{proof}
\begin{definition}[Matches Given Absorption Wavelength Readouts]
  \label{def:physics:phyx-mini-0587:phyxminiproblems-problemphyxmini0587-matchesgivenabsorptionwavelengthreadouts}
  \lean{PhyXMiniProblems.ProblemPhyXMini0587.MatchesGivenAbsorptionWavelengthReadouts}
  \uses{def:physics:phyx-mini-0587:phyxminiproblems-problemphyxmini0587-infinitewellabsorptionsetup, def:physics:phyx-mini-0587:phyxminiproblems-problemphyxmini0587-agreeswithtwodecimalnanometerreadout}
  The five wavelength values supplied in the problem statement
\end{definition}
\begin{proof}
  This is the declared model definition of Matches Given Absorption Wavelength Readouts.
\end{proof}
\begin{definition}[Matches Supplied Absorption Figure]
  \label{def:physics:phyx-mini-0587:phyxminiproblems-problemphyxmini0587-matchessuppliedabsorptionfigure}
  \lean{PhyXMiniProblems.ProblemPhyXMini0587.MatchesSuppliedAbsorptionFigure}
  \uses{def:physics:phyx-mini-0587:phyxminiproblems-problemphyxmini0587-infinitewellabsorptionfigure}
  Primary-image evidence from 587.png: a wavelength axis in nanometres with a zero tick, subdivision marks, and five red labelled points ordered lambdaE, lambdaD, lambdaC, lambdaB, lambdaA from left to right
\end{definition}
\begin{proof}
  This is the declared model definition of Matches Supplied Absorption Figure.
\end{proof}
\begin{definition}[Matches Five Longest Absorption Transitions]
  \label{def:physics:phyx-mini-0587:phyxminiproblems-problemphyxmini0587-matchesfivelongestabsorptiontransitions}
  \lean{PhyXMiniProblems.ProblemPhyXMini0587.MatchesFiveLongestAbsorptionTransitions}
  \uses{def:physics:phyx-mini-0587:phyxminiproblems-problemphyxmini0587-infinitewellabsorptionsetup}
  The five longest upward absorptions from n = 2 reach the next five energy levels. This identifies the physical transition represented by each plotted label without assuming anything about the requested well width
\end{definition}
\begin{proof}
  This is the declared model definition of Matches Five Longest Absorption Transitions.
\end{proof}
\begin{definition}[Uses Standard Reference Data]
  \label{def:physics:phyx-mini-0587:phyxminiproblems-problemphyxmini0587-usesstandardreferencedata}
  \lean{PhyXMiniProblems.ProblemPhyXMini0587.UsesStandardReferenceData}
  \uses{def:physics:phyx-mini-0587:phyxminiproblems-problemphyxmini0587-massinkilograms, def:physics:phyx-mini-0587:phyxminiproblems-problemphyxmini0587-actioninjouleseconds, def:physics:phyx-mini-0587:phyxminiproblems-problemphyxmini0587-speedinmeterspersecond, def:physics:phyx-mini-0587:phyxminiproblems-problemphyxmini0587-infinitewellabsorptionsetup}
  Electron mass, reduced Planck constant, and vacuum light-speed calibrations
\end{definition}
\begin{proof}
  This is the declared model definition of Uses Standard Reference Data.
\end{proof}
\begin{definition}[Has Physical Infinite Well Parameters]
  \label{def:physics:phyx-mini-0587:phyxminiproblems-problemphyxmini0587-hasphysicalinfinitewellparameters}
  \lean{PhyXMiniProblems.ProblemPhyXMini0587.HasPhysicalInfiniteWellParameters}
  \uses{def:physics:phyx-mini-0587:phyxminiproblems-problemphyxmini0587-lengthinmeters, def:physics:phyx-mini-0587:phyxminiproblems-problemphyxmini0587-massinkilograms, def:physics:phyx-mini-0587:phyxminiproblems-problemphyxmini0587-actioninjouleseconds, def:physics:phyx-mini-0587:phyxminiproblems-problemphyxmini0587-speedinmeterspersecond, def:physics:phyx-mini-0587:phyxminiproblems-problemphyxmini0587-infinitewellabsorptionsetup}
  Positivity conditions selecting a nondegenerate physical experiment
\end{definition}
\begin{proof}
  This is the declared model definition of Has Physical Infinite Well Parameters.
\end{proof}
\begin{definition}[Satisfies Infinite Square Well Spectrum]
  \label{def:physics:phyx-mini-0587:phyxminiproblems-problemphyxmini0587-satisfiesinfinitesquarewellspectrum}
  \lean{PhyXMiniProblems.ProblemPhyXMini0587.SatisfiesInfiniteSquareWellSpectrum}
  \uses{def:physics:phyx-mini-0587:phyxminiproblems-problemphyxmini0587-lengthinmeters, def:physics:phyx-mini-0587:phyxminiproblems-problemphyxmini0587-massinkilograms, def:physics:phyx-mini-0587:phyxminiproblems-problemphyxmini0587-actioninjouleseconds, def:physics:phyx-mini-0587:phyxminiproblems-problemphyxmini0587-energyinjoules, def:physics:phyx-mini-0587:phyxminiproblems-problemphyxmini0587-infinitewellabsorptionsetup}
  The one-dimensional infinite-square-well spectrum E\_n = n\textasciicircum{}2 * pi\textasciicircum{}2 * hbar\textasciicircum{}2 / (2 m L\textasciicircum{}2) for positive quantum numbers. Physlib currently has no dedicated infinite- well spectrum declaration, so this is a local governing-law interface
\end{definition}
\begin{proof}
  This is the declared model definition of Satisfies Infinite Square Well Spectrum.
\end{proof}
\begin{definition}[Satisfies Single Photon Absorption Law]
  \label{def:physics:phyx-mini-0587:phyxminiproblems-problemphyxmini0587-satisfiessinglephotonabsorptionlaw}
  \lean{PhyXMiniProblems.ProblemPhyXMini0587.SatisfiesSinglePhotonAbsorptionLaw}
  \uses{def:physics:phyx-mini-0587:phyxminiproblems-problemphyxmini0587-lengthinmeters, def:physics:phyx-mini-0587:phyxminiproblems-problemphyxmini0587-actioninjouleseconds, def:physics:phyx-mini-0587:phyxminiproblems-problemphyxmini0587-speedinmeterspersecond, def:physics:phyx-mini-0587:phyxminiproblems-problemphyxmini0587-energyinjoules, def:physics:phyx-mini-0587:phyxminiproblems-problemphyxmini0587-infinitewellabsorptionsetup}
  Each plotted line is a one-photon absorption: the photon energy is 2 * pi * hbar * c / lambda and equals the electron's final-minus-initial level energy. This is a governing law, not the requested numerical answer
\end{definition}
\begin{proof}
  This is the declared model definition of Satisfies Single Photon Absorption Law.
\end{proof}
\begin{lemma}[well Width Squared from longest Absorption Line]
  \label{lem:physics:phyx-mini-0587:phyxminiproblems-problemphyxmini0587-wellwidthsquared-from-longestabsorptionline}
  \lean{PhyXMiniProblems.ProblemPhyXMini0587.wellWidthSquared_from_longestAbsorptionLine}
  \uses{def:physics:phyx-mini-0587:phyxminiproblems-problemphyxmini0587-lengthinmeters, def:physics:phyx-mini-0587:phyxminiproblems-problemphyxmini0587-massinkilograms, def:physics:phyx-mini-0587:phyxminiproblems-problemphyxmini0587-actioninjouleseconds, def:physics:phyx-mini-0587:phyxminiproblems-problemphyxmini0587-speedinmeterspersecond, def:physics:phyx-mini-0587:phyxminiproblems-problemphyxmini0587-infinitewellabsorptionsetup, def:physics:phyx-mini-0587:phyxminiproblems-problemphyxmini0587-matchesfirstexcitedelectroninfinitewell, def:physics:phyx-mini-0587:phyxminiproblems-problemphyxmini0587-matchesfivelongestabsorptiontransitions, def:physics:phyx-mini-0587:phyxminiproblems-problemphyxmini0587-hasphysicalinfinitewellparameters, def:physics:phyx-mini-0587:phyxminiproblems-problemphyxmini0587-satisfiesinfinitesquarewellspectrum, def:physics:phyx-mini-0587:phyxminiproblems-problemphyxmini0587-satisfiessinglephotonabsorptionlaw}
  The n = 2 to n = 3 line implies the coherent-SI relation L\textasciicircum{}2 = 5 * pi * hbar * lambdaA / (4 m c). This is a derived conclusion and does not contain the displayed 350 pm answer
\end{lemma}
\begin{proof}
  The conclusion follows from the governing relations and figure readouts named in the statement of well Width Squared from longest Absorption Line.
\end{proof}
\begin{definition}[Answer Choice]
  \label{def:physics:phyx-mini-0587:phyxminiproblems-problemphyxmini0587-answerchoice}
  \lean{PhyXMiniProblems.ProblemPhyXMini0587.AnswerChoice}
  Labels of the four well-width choices printed in the problem
\end{definition}
\begin{proof}
  This is the declared model definition of Answer Choice.
\end{proof}
\begin{definition}[displayed Well Width Picometers]
  \label{def:physics:phyx-mini-0587:phyxminiproblems-problemphyxmini0587-displayedwellwidthpicometers}
  \lean{PhyXMiniProblems.ProblemPhyXMini0587.displayedWellWidthPicometers}
  \uses{def:physics:phyx-mini-0587:phyxminiproblems-problemphyxmini0587-answerchoice}
  The well width printed beside each answer label, measured in picometres
\end{definition}
\begin{proof}
  This is the declared model definition of displayed Well Width Picometers.
\end{proof}
\begin{definition}[recorded Dataset Answer]
  \label{def:physics:phyx-mini-0587:phyxminiproblems-problemphyxmini0587-recordeddatasetanswer}
  \lean{PhyXMiniProblems.ProblemPhyXMini0587.recordedDatasetAnswer}
  \uses{def:physics:phyx-mini-0587:phyxminiproblems-problemphyxmini0587-answerchoice}
  Dataset metadata recording answer C; deliberately not used as a premise
\end{definition}
\begin{proof}
  This is the declared model definition of recorded Dataset Answer.
\end{proof}
\begin{definition}[Agrees When Rounded To Nearest Picometer]
  \label{def:physics:phyx-mini-0587:phyxminiproblems-problemphyxmini0587-agreeswhenroundedtonearestpicometer}
  \lean{PhyXMiniProblems.ProblemPhyXMini0587.AgreesWhenRoundedToNearestPicometer}
  \uses{def:physics:phyx-mini-0587:phyxminiproblems-problemphyxmini0587-lengthinpicometers, def:physics:phyx-mini-0587:phyxminiproblems-problemphyxmini0587-infinitewellabsorptionsetup, def:physics:phyx-mini-0587:phyxminiproblems-problemphyxmini0587-answerchoice, def:physics:phyx-mini-0587:phyxminiproblems-problemphyxmini0587-displayedwellwidthpicometers}
  Agreement with a well width displayed to the nearest whole picometre
\end{definition}
\begin{proof}
  This is the declared model definition of Agrees When Rounded To Nearest Picometer.
\end{proof}
\begin{definition}[Is Unique Nearest Well Width Choice]
  \label{def:physics:phyx-mini-0587:phyxminiproblems-problemphyxmini0587-isuniquenearestwellwidthchoice}
  \lean{PhyXMiniProblems.ProblemPhyXMini0587.IsUniqueNearestWellWidthChoice}
  \uses{def:physics:phyx-mini-0587:phyxminiproblems-problemphyxmini0587-lengthinpicometers, def:physics:phyx-mini-0587:phyxminiproblems-problemphyxmini0587-infinitewellabsorptionsetup, def:physics:phyx-mini-0587:phyxminiproblems-problemphyxmini0587-answerchoice, def:physics:phyx-mini-0587:phyxminiproblems-problemphyxmini0587-displayedwellwidthpicometers}
  The selected answer is strictly nearer than every other displayed width
\end{definition}
\begin{proof}
  This is the declared model definition of Is Unique Nearest Well Width Choice.
\end{proof}
% --- Archon declaration topology end ---
% --- Archon physics formalization source end ---
